% =============================================================================
% Preambles/paper-header.tex — shared LaTeX preamble for MANUSCRIPTS (not slides)
%
% Sibling to Preambles/header.tex (the Beamer/Quarto slide preamble). Use this
% one for a research paper via:
%     % =============================================================================
% Preambles/paper-header.tex — shared LaTeX preamble for MANUSCRIPTS (not slides)
%
% Sibling to Preambles/header.tex (the Beamer/Quarto slide preamble). Use this
% one for a research paper via:
%     % =============================================================================
% Preambles/paper-header.tex — shared LaTeX preamble for MANUSCRIPTS (not slides)
%
% Sibling to Preambles/header.tex (the Beamer/Quarto slide preamble). Use this
% one for a research paper via:
%     % =============================================================================
% Preambles/paper-header.tex — shared LaTeX preamble for MANUSCRIPTS (not slides)
%
% Sibling to Preambles/header.tex (the Beamer/Quarto slide preamble). Use this
% one for a research paper via:
%     \input{../Preambles/paper-header}
% from templates/paper/paper-template.tex or any manuscript in Paper/.
%
% DESIGN GOAL: a professional working-paper look out of the box — no visible
% hyperlink boxes, a real title/author grid, booktabs tables, natbib citations.
% See .claude/rules/paper-writing-craft.md §4 for the rationale behind each
% choice below.
% =============================================================================

\usepackage{iftex}
\ifPDFTeX
  \usepackage[utf8]{inputenc}
  \usepackage[T1]{fontenc}
\fi

\usepackage[english]{babel}
\usepackage{amsmath, amssymb, amsthm}
\usepackage{booktabs}      % \toprule/\midrule/\bottomrule — never hand-drawn \hline grids
\usepackage{graphicx}
\usepackage{geometry}
\geometry{margin=1in}

% -----------------------------------------------------------------------------
% Font — Latin Modern by default (a cleaner-rendering Computer Modern, not a
% jarring departure). Uncomment ONE warmer-serif alternative if you prefer it;
% keep both preamble and any Quarto-side font choice consistent if you also
% maintain a Quarto mirror of this content.
% -----------------------------------------------------------------------------
\usepackage{lmodern}
% \usepackage{mathpazo}                    % Palatino-based text + math
% \usepackage{newtxtext,newtxmath}         % Times-based text + math

% -----------------------------------------------------------------------------
% Epigraph support — optional narrative device (paper-writing-craft.md §2).
% Use sparingly: \epigraph{quote}{attribution} once, in the introduction.
% -----------------------------------------------------------------------------
\usepackage{epigraph}
\setlength{\epigraphwidth}{0.85\textwidth}
\renewcommand{\epigraphflush}{flushleft}
\renewcommand{\epigraphrule}{0pt}

% -----------------------------------------------------------------------------
% Hyperref — loaded near the end per standard practice. `hidelinks` suppresses
% BOTH the colored text and the visible frame that hyperref's un-configured
% default draws around every citation and cross-reference. That boxed-link
% look is the single most common "unpolished draft" tell — see
% paper-writing-craft.md §4. Swap to colorlinks=true below if you prefer subtle
% color over fully plain text; never ship the bare, unconfigured default.
% -----------------------------------------------------------------------------
\usepackage[hidelinks]{hyperref}
% \usepackage{hyperref}
% \hypersetup{colorlinks=true, linkcolor=black, citecolor=black, urlcolor=black}

\usepackage[authoryear,round]{natbib}   % author-year citations; swap `round` for `square` to taste
\bibliographystyle{apalike}             % or econometrica.bst / aer.bst if your journal target ships one

% -----------------------------------------------------------------------------
% Title / author / affiliation block — a centered multi-column grid rather
% than a single flush-left author line, for papers with more than one author.
% Usage in the manuscript:
%
%   \title{[TITLE]}
%   \authorgrid{
%     \shortstack{[Author One] \\ \small [Affiliation One]} &
%     \shortstack{[Author Two] \\ \small [Affiliation Two]}
%   }
%   \date{[DATE]}
%   \maketitle
%   \thispagestyle{empty}
% -----------------------------------------------------------------------------
\newcommand{\authorgrid}[1]{%
  \author{\begin{tabular}[t]{@{}c@{\hspace{2.5em}}c@{}}#1\end{tabular}}%
}

% -----------------------------------------------------------------------------
% Abstract environment — the standard `article` class's own `abstract`
% environment already renders a centered bold "Abstract" header followed by an
% indented, distinguishable block, which is exactly the professional
% convention paper-writing-craft.md §4 asks for. No redefinition needed —
% just use \begin{abstract}...\end{abstract} in the manuscript body.
% -----------------------------------------------------------------------------

% -----------------------------------------------------------------------------
% Acknowledgment footnote convention — attach to \title via \thanks{...}, e.g.
%   \title{[TITLE]\thanks{We thank [names] for helpful comments. [RA credits].}}
% renders as a numbered footnote on the first page, not inline text.
% -----------------------------------------------------------------------------

% from templates/paper/paper-template.tex or any manuscript in Paper/.
%
% DESIGN GOAL: a professional working-paper look out of the box — no visible
% hyperlink boxes, a real title/author grid, booktabs tables, natbib citations.
% See .claude/rules/paper-writing-craft.md §4 for the rationale behind each
% choice below.
% =============================================================================

\usepackage{iftex}
\ifPDFTeX
  \usepackage[utf8]{inputenc}
  \usepackage[T1]{fontenc}
\fi

\usepackage[english]{babel}
\usepackage{amsmath, amssymb, amsthm}
\usepackage{booktabs}      % \toprule/\midrule/\bottomrule — never hand-drawn \hline grids
\usepackage{graphicx}
\usepackage{geometry}
\geometry{margin=1in}

% -----------------------------------------------------------------------------
% Font — Latin Modern by default (a cleaner-rendering Computer Modern, not a
% jarring departure). Uncomment ONE warmer-serif alternative if you prefer it;
% keep both preamble and any Quarto-side font choice consistent if you also
% maintain a Quarto mirror of this content.
% -----------------------------------------------------------------------------
\usepackage{lmodern}
% \usepackage{mathpazo}                    % Palatino-based text + math
% \usepackage{newtxtext,newtxmath}         % Times-based text + math

% -----------------------------------------------------------------------------
% Epigraph support — optional narrative device (paper-writing-craft.md §2).
% Use sparingly: \epigraph{quote}{attribution} once, in the introduction.
% -----------------------------------------------------------------------------
\usepackage{epigraph}
\setlength{\epigraphwidth}{0.85\textwidth}
\renewcommand{\epigraphflush}{flushleft}
\renewcommand{\epigraphrule}{0pt}

% -----------------------------------------------------------------------------
% Hyperref — loaded near the end per standard practice. `hidelinks` suppresses
% BOTH the colored text and the visible frame that hyperref's un-configured
% default draws around every citation and cross-reference. That boxed-link
% look is the single most common "unpolished draft" tell — see
% paper-writing-craft.md §4. Swap to colorlinks=true below if you prefer subtle
% color over fully plain text; never ship the bare, unconfigured default.
% -----------------------------------------------------------------------------
\usepackage[hidelinks]{hyperref}
% \usepackage{hyperref}
% \hypersetup{colorlinks=true, linkcolor=black, citecolor=black, urlcolor=black}

\usepackage[authoryear,round]{natbib}   % author-year citations; swap `round` for `square` to taste
\bibliographystyle{apalike}             % or econometrica.bst / aer.bst if your journal target ships one

% -----------------------------------------------------------------------------
% Title / author / affiliation block — a centered multi-column grid rather
% than a single flush-left author line, for papers with more than one author.
% Usage in the manuscript:
%
%   \title{[TITLE]}
%   \authorgrid{
%     \shortstack{[Author One] \\ \small [Affiliation One]} &
%     \shortstack{[Author Two] \\ \small [Affiliation Two]}
%   }
%   \date{[DATE]}
%   \maketitle
%   \thispagestyle{empty}
% -----------------------------------------------------------------------------
\newcommand{\authorgrid}[1]{%
  \author{\begin{tabular}[t]{@{}c@{\hspace{2.5em}}c@{}}#1\end{tabular}}%
}

% -----------------------------------------------------------------------------
% Abstract environment — the standard `article` class's own `abstract`
% environment already renders a centered bold "Abstract" header followed by an
% indented, distinguishable block, which is exactly the professional
% convention paper-writing-craft.md §4 asks for. No redefinition needed —
% just use \begin{abstract}...\end{abstract} in the manuscript body.
% -----------------------------------------------------------------------------

% -----------------------------------------------------------------------------
% Acknowledgment footnote convention — attach to \title via \thanks{...}, e.g.
%   \title{[TITLE]\thanks{We thank [names] for helpful comments. [RA credits].}}
% renders as a numbered footnote on the first page, not inline text.
% -----------------------------------------------------------------------------

% from templates/paper/paper-template.tex or any manuscript in Paper/.
%
% DESIGN GOAL: a professional working-paper look out of the box — no visible
% hyperlink boxes, a real title/author grid, booktabs tables, natbib citations.
% See .claude/rules/paper-writing-craft.md §4 for the rationale behind each
% choice below.
% =============================================================================

\usepackage{iftex}
\ifPDFTeX
  \usepackage[utf8]{inputenc}
  \usepackage[T1]{fontenc}
\fi

\usepackage[english]{babel}
\usepackage{amsmath, amssymb, amsthm}
\usepackage{booktabs}      % \toprule/\midrule/\bottomrule — never hand-drawn \hline grids
\usepackage{graphicx}
\usepackage{geometry}
\geometry{margin=1in}

% -----------------------------------------------------------------------------
% Font — Latin Modern by default (a cleaner-rendering Computer Modern, not a
% jarring departure). Uncomment ONE warmer-serif alternative if you prefer it;
% keep both preamble and any Quarto-side font choice consistent if you also
% maintain a Quarto mirror of this content.
% -----------------------------------------------------------------------------
\usepackage{lmodern}
% \usepackage{mathpazo}                    % Palatino-based text + math
% \usepackage{newtxtext,newtxmath}         % Times-based text + math

% -----------------------------------------------------------------------------
% Epigraph support — optional narrative device (paper-writing-craft.md §2).
% Use sparingly: \epigraph{quote}{attribution} once, in the introduction.
% -----------------------------------------------------------------------------
\usepackage{epigraph}
\setlength{\epigraphwidth}{0.85\textwidth}
\renewcommand{\epigraphflush}{flushleft}
\renewcommand{\epigraphrule}{0pt}

% -----------------------------------------------------------------------------
% Hyperref — loaded near the end per standard practice. `hidelinks` suppresses
% BOTH the colored text and the visible frame that hyperref's un-configured
% default draws around every citation and cross-reference. That boxed-link
% look is the single most common "unpolished draft" tell — see
% paper-writing-craft.md §4. Swap to colorlinks=true below if you prefer subtle
% color over fully plain text; never ship the bare, unconfigured default.
% -----------------------------------------------------------------------------
\usepackage[hidelinks]{hyperref}
% \usepackage{hyperref}
% \hypersetup{colorlinks=true, linkcolor=black, citecolor=black, urlcolor=black}

\usepackage[authoryear,round]{natbib}   % author-year citations; swap `round` for `square` to taste
\bibliographystyle{apalike}             % or econometrica.bst / aer.bst if your journal target ships one

% -----------------------------------------------------------------------------
% Title / author / affiliation block — a centered multi-column grid rather
% than a single flush-left author line, for papers with more than one author.
% Usage in the manuscript:
%
%   \title{[TITLE]}
%   \authorgrid{
%     \shortstack{[Author One] \\ \small [Affiliation One]} &
%     \shortstack{[Author Two] \\ \small [Affiliation Two]}
%   }
%   \date{[DATE]}
%   \maketitle
%   \thispagestyle{empty}
% -----------------------------------------------------------------------------
\newcommand{\authorgrid}[1]{%
  \author{\begin{tabular}[t]{@{}c@{\hspace{2.5em}}c@{}}#1\end{tabular}}%
}

% -----------------------------------------------------------------------------
% Abstract environment — the standard `article` class's own `abstract`
% environment already renders a centered bold "Abstract" header followed by an
% indented, distinguishable block, which is exactly the professional
% convention paper-writing-craft.md §4 asks for. No redefinition needed —
% just use \begin{abstract}...\end{abstract} in the manuscript body.
% -----------------------------------------------------------------------------

% -----------------------------------------------------------------------------
% Acknowledgment footnote convention — attach to \title via \thanks{...}, e.g.
%   \title{[TITLE]\thanks{We thank [names] for helpful comments. [RA credits].}}
% renders as a numbered footnote on the first page, not inline text.
% -----------------------------------------------------------------------------

% from templates/paper/paper-template.tex or any manuscript in Paper/.
%
% DESIGN GOAL: a professional working-paper look out of the box — no visible
% hyperlink boxes, a real title/author grid, booktabs tables, natbib citations.
% See .claude/rules/paper-writing-craft.md §4 for the rationale behind each
% choice below.
% =============================================================================

\usepackage{iftex}
\ifPDFTeX
  \usepackage[utf8]{inputenc}
  \usepackage[T1]{fontenc}
\fi

\usepackage[english]{babel}
\usepackage{amsmath, amssymb, amsthm}
\usepackage{booktabs}      % \toprule/\midrule/\bottomrule — never hand-drawn \hline grids
\usepackage{graphicx}
\usepackage{geometry}
\geometry{margin=1in}

% -----------------------------------------------------------------------------
% Font — Latin Modern by default (a cleaner-rendering Computer Modern, not a
% jarring departure). Uncomment ONE warmer-serif alternative if you prefer it;
% keep both preamble and any Quarto-side font choice consistent if you also
% maintain a Quarto mirror of this content.
% -----------------------------------------------------------------------------
\usepackage{lmodern}
% \usepackage{mathpazo}                    % Palatino-based text + math
% \usepackage{newtxtext,newtxmath}         % Times-based text + math

% -----------------------------------------------------------------------------
% Epigraph support — optional narrative device (paper-writing-craft.md §2).
% Use sparingly: \epigraph{quote}{attribution} once, in the introduction.
% -----------------------------------------------------------------------------
\usepackage{epigraph}
\setlength{\epigraphwidth}{0.85\textwidth}
\renewcommand{\epigraphflush}{flushleft}
\renewcommand{\epigraphrule}{0pt}

% -----------------------------------------------------------------------------
% Hyperref — loaded near the end per standard practice. `hidelinks` suppresses
% BOTH the colored text and the visible frame that hyperref's un-configured
% default draws around every citation and cross-reference. That boxed-link
% look is the single most common "unpolished draft" tell — see
% paper-writing-craft.md §4. Swap to colorlinks=true below if you prefer subtle
% color over fully plain text; never ship the bare, unconfigured default.
% -----------------------------------------------------------------------------
\usepackage[hidelinks]{hyperref}
% \usepackage{hyperref}
% \hypersetup{colorlinks=true, linkcolor=black, citecolor=black, urlcolor=black}

\usepackage[authoryear,round]{natbib}   % author-year citations; swap `round` for `square` to taste
\bibliographystyle{apalike}             % or econometrica.bst / aer.bst if your journal target ships one

% -----------------------------------------------------------------------------
% Title / author / affiliation block — a centered multi-column grid rather
% than a single flush-left author line, for papers with more than one author.
% Usage in the manuscript:
%
%   \title{[TITLE]}
%   \authorgrid{
%     \shortstack{[Author One] \\ \small [Affiliation One]} &
%     \shortstack{[Author Two] \\ \small [Affiliation Two]}
%   }
%   \date{[DATE]}
%   \maketitle
%   \thispagestyle{empty}
% -----------------------------------------------------------------------------
\newcommand{\authorgrid}[1]{%
  \author{\begin{tabular}[t]{@{}c@{\hspace{2.5em}}c@{}}#1\end{tabular}}%
}

% -----------------------------------------------------------------------------
% Abstract environment — the standard `article` class's own `abstract`
% environment already renders a centered bold "Abstract" header followed by an
% indented, distinguishable block, which is exactly the professional
% convention paper-writing-craft.md §4 asks for. No redefinition needed —
% just use \begin{abstract}...\end{abstract} in the manuscript body.
% -----------------------------------------------------------------------------

% -----------------------------------------------------------------------------
% Acknowledgment footnote convention — attach to \title via \thanks{...}, e.g.
%   \title{[TITLE]\thanks{We thank [names] for helpful comments. [RA credits].}}
% renders as a numbered footnote on the first page, not inline text.
% -----------------------------------------------------------------------------
